\article{约化对偶对的例子}{Examples of Dual Reductive Pairs}{Stephen Gelbart}{1}{287--296}
\noindent\textit{作者附注：Alfred P. Sloan 研究员。}

本文以具体、直观的例子介绍约化对偶对（dual reductive pair），以及由它们产生的自守形式（automorphic form）和群表示（group representation），作为 Howe 论文 [Ho 1] 的补充。感谢 Howe 向我讲解他的振子表示（oscillator representation）理论。

\paragraph*{本文目录.}
1. 历史评述（原第287页）；2. Schrödinger 模型（288页）；3. 对 $(\operatorname{Sp}(V_0),O(U_0))$（289页）；4. 阿代尔化（290页）；5. 局部例子（291页）；6. 整体例子（293页）。

\origsection{1}{历史评述（Historical remarks）}
1964 年，Weil 为了给经典 theta 函数
\[
 \theta(z)=\sum_{n=-\infty}^\infty\exp(2\pi i n^2z)
\]
奠定群论基础，引入了辛群（symplectic group）的一个自然表示，即 Segal--Shale--Weil--metaplectic 振子表示。对于 $\operatorname{SL}(2,F)$，这个表示由以下公式给出：
\[
 T_\chi\!\begin{pmatrix}1&b\\0&1\end{pmatrix}\Phi(x)
 =\chi(\tfrac12 bx^2)\Phi(x)
\]
以及
\[
 T_\chi\!\begin{pmatrix}0&1\\-1&0\end{pmatrix}\Phi(x)
 =\gamma\widehat\Phi(x)
\]
作用在 $L^2(F)$ 中。这里 $\chi$ 是 $F$ 的一个固定非平凡特征（character），$\gamma$ 是依赖于 $\chi$ 的一个八次单位根。对于 $\operatorname{Sp}(2m,\R)$，振子表示通过类似的算子作用于 $L^2(\R^m)$。

1966 年，Shalika 和 Tanaka 将 $T_\chi$ 与 $T_{\chi'}$ 作张量积，从而在 $L^2(F^2)$ 中得到 $\operatorname{SL}(2,F)$ 的一个表示；这个表示与 $O(2)$ 的自然作用对易，其分解可由后者的表示论来描述。这样，Shalika 和 Tanaka 构造出 $\operatorname{SL}(2,F)$ 的离散系列表示（discrete series representation），并重新表述了 Hecke 和 Maass 关于带广义 Hecke 特征（Grössencharakter）的 theta 级数构造的早期结果。对于 $\operatorname{Sp}(2m,\R)$，类似的 $2m$ 重振子表示张量积通过算子
\begin{equation}
\begin{aligned}
 \omega_\chi\!\begin{pmatrix}I_m&N\\0&I_m\end{pmatrix}\Phi(X)
 &=\exp(\pi i\operatorname{tr}XNX^t)\Phi(X),\\
 \omega_\chi\!\begin{pmatrix}A&0\\0&(A^t)^{-1}\end{pmatrix}\Phi(X)
 &=|\det A|^m\Phi(XA)
\end{aligned}\tag{1.1}
\end{equation}
以及
\begin{equation}
 \omega_\chi\!\begin{pmatrix}0&I_m\\-I_m&0\end{pmatrix}\Phi(X)
 =i^{-m^2}\widehat\Phi(X)\tag{1.2}
\end{equation}
作用于 $L^2(\R^{2m\times m})$。这个表示分解为 $\operatorname{Sp}(2m,\R)$ 的离散系列表示，由 $O(2m,\R)$ 的某些不可约表示（irreducible representation）索引，具体而言，$O(2m,\R)$ 通过矩阵左乘自然地作用于 $L^2(\R^{2m\times m})$，指标正是这一作用中出现的不可约表示；参见 [Ge 1] 和 [Sa]。

Howe 最近发展出的振子表示理论重新表述并精化了 Weil 的理论，从而为上述结果提供统一解释，并指出进一步研究的方向。约化对偶对是这一理论的核心概念之一，它将上述例子中辛群与正交群之间的自然对偶性表述为一般的数学结构。要更具体地描述这一现象，Schrödinger 模型（Schrödinger model）的概念至关重要。

\origsection{2}{Schrödinger 模型（Schrödinger models）}
我们采用 [Ho 1] 的记号和定义。其中，$F$ 是特征不为 $2$ 的局部域（local field），$V$ 是定义在 $F$ 上的辛空间，$H(V)$ 是与 $V$ 相联系的 Heisenberg 群（Heisenberg group），$\chi$ 是 $F$ 的一个固定非平凡特征，$\rho_\chi$ 是与 $\chi$ 相联系的 $H$ 的不可约酉表示（unitary representation），$\operatorname{Sp}(V)$ 是辛型 $\langle\ ,\ \rangle$ 的等距群（isometry group）。回忆，$H(V)$ 是以 $F$ 为中心的 $V$ 的中心扩张（central extension）。

令 $\overline{\operatorname{Sp}}$ 表示 $\operatorname{Sp}(V)$ 的非平凡二重覆盖（double cover）；其存在性和唯一性见 [Mo]。\footnote{\textbf{译者注（本次校订）：}这里应排除 $F=\C$。复辛群 $\operatorname{Sp}(2n,\C)$ 单连通，其 metaplectic 二重覆盖分裂；此时振子表示可直接视为辛群的表示。前篇第2节已单独说明这一例外。}振子表示 $\omega_\chi$ 是 $\overline{\operatorname{Sp}}$ 在 $\rho_\chi$ 的空间中的酉表示，满足：对所有 $g\in\overline{\operatorname{Sp}}$ 和 $h\in H$，
\begin{equation}
 \omega_\chi(g)\rho_\chi(h)\omega_\chi(g^{-1})=\rho_\chi(g(h)).\tag{2.1}
\end{equation}
这一表示在酉等价意义下唯一，而且是“真正的”（genuine），即不能下降为 $\operatorname{Sp}(V)$ 的表示。相应的 $\operatorname{Sp}(V)$ 的乘子表示（multiplier representation）仍记为 $\omega_\chi$。

要描述 $\omega_\chi$（及 $\rho_\chi$）的 Schrödinger 模型，需要 $V$ 的完全极化（complete polarization）这一概念。完全极化是满足下列条件的一对子空间 $(X,Y)$：
\begin{enumerate}
\item[(i)] $X,Y$ 对 $\langle\ ,\ \rangle$ 各向同性（isotropic），即 $\langle\ ,\ \rangle$ 在 $X$ 和 $Y$ 上均平凡；
\item[(ii)] $X\oplus Y=V$。
\end{enumerate}
选定这一极化后，$\rho_\chi$ 可实现于 $L^2(X)$，$H(V)$ 的作用为\footnote{\textbf{译者注（本次校订）：}式 (2.2) 第二行的条件应为 $y\in Y$，而不是 $y\in V$。平移作用与乘法作用分别对应所选极化的两个子空间 $X$ 和 $Y$。}
\begin{equation}
\begin{aligned}
 \rho_\chi((x,0))f(x')&=f(x'-x),&&x,x'\in X,\\
 \rho_\chi((y,0))f(x')&=\chi(\langle y,x'\rangle)f(x'),&&y\in V,\ x'\in X,\\
 \rho_\chi((0,t))f(x')&=\chi(t)f(x'),&&t\in F,\ x'\in X.
\end{aligned}\tag{2.2}
\end{equation}
相应的 $\omega_\chi$ 的 Schrödinger 模型较为复杂，我们先考察它在 $\operatorname{Sp}(V)$ 的几个特殊子群上的限制。

令 $P(Y)$（分别为 $P(X)$）表示 $\operatorname{Sp}(V)$ 中保持 $Y$（分别为 $X$）的子群，令 $N(Y)$ 表示在 $Y$ 上作用平凡的子群。若
\[
 M=M(X,Y)=P(X)\cap P(Y),
\]
则 $P(Y)$ 是 $M$ 与 $N(Y)$ 的半直积（semidirect product）。利用 (2.1) 和 (2.2)，可得
\begin{equation}
 \omega_\chi(n)f(x)=\chi\bigl(-\tfrac12\langle x,nx\rangle\bigr)f(x),
 \qquad x\in X,\ n\in N(Y),\tag{2.3}
\end{equation}
以及
\begin{equation}
 \omega_\chi(m)f(x)=\mu(m)f(m^{-1}x),\qquad x\in X,\ m\in M.\tag{2.4}
\end{equation}
因此，在 $\omega_\chi$ 的 Schrödinger 模型中，$N$ 通过乘以一个二次特征作用，而 $M$ 的作用由线性换元及一个标量因子给出；因子 $\mu(m)$ 大致是 $m$ 在 $X$ 上作用的模（modulus）的平方根，参见 (1.1)。

对于 $P(Y)$ 之外的元素，$\omega_\chi$ 的作用可参看后文的具体例子。回忆，形如 $\left(\begin{smallmatrix}1&b\\0&1\end{smallmatrix}\right)$ 和 $\left(\begin{smallmatrix}0&1\\-1&0\end{smallmatrix}\right)$ 的矩阵生成 $\operatorname{SL}(2,F)$。

\origsection{3}{对 $(\operatorname{Sp}(V_0),O(U_0))$}
约化对偶对的一般概念见 [Ho 1]。本文的基本例子将上面 $\operatorname{Sp}(2m,\R)$ 的特殊性质概括为一个统一的构造。

设 $(V_0,\langle\ ,\ \rangle_0)$ 为辛向量空间，$(U_0,(\ ,\ ))$ 为正交向量空间。令 $V=V_0\otimes U_0$，并通过
\[
 \langle v\otimes u,v'\otimes u'\rangle
 =\langle v,v'\rangle_0\langle u,u'\rangle_0,
 \qquad v,v'\in V_0,\ u,u'\in U_0,
\]
在 $V$ 上定义双线性型 $\langle\ ,\ \rangle$。\footnote{\textbf{译者注（本次校订）：}右端第二个因子应为先前定义的对称双线性型 $(u,u^{\prime})$，即 $\langle v\otimes u,v^{\prime}\otimes u^{\prime}\rangle=\langle v,v^{\prime}\rangle_0(u,u^{\prime})$。一个交替型与一个对称型的张量积才给出所需的交替型。}那么 $(V,\langle\ ,\ \rangle)$ 是辛空间，$(\operatorname{Sp}(V_0),O(U_0))$ 是 $\operatorname{Sp}(V)$ 中自然的约化对偶对。特别地，群 $\operatorname{Sp}(V_0)$ 与 $O(U_0)$ 在 $\operatorname{Sp}(V)$ 中互为中心化子（centralizer）。$\operatorname{Sp}(V)$ 中的 $(\operatorname{Sp}(2m,\R),O(2m,\R))$ 就是这样一对，其中 $V=\R^{2m}\otimes\R^{2m}=V_0\otimes U_0$。

若 $(X_0,Y_0)$ 是 $V_0$ 的完全极化，令 $X=X_0\otimes U_0$，$Y=Y_0\otimes U_0$。那么 $(X,Y)$ 是 $V$ 的完全极化。此外，由第2节可知，$\omega_\chi$ 在 $L^2(X)$ 中的 Schrödinger 模型使 $O(U_0)$ 作线性作用。另一方面，若 $(X_0,Y_0)$ 是 $U_0$ 的一个完全极化（这只有在 $U_0$ 分裂（split）时才可能），则空间 $X=V_0\otimes X_0$ 和 $Y=V_0\otimes Y_0$ 给出 $V$ 的一个极化，在相应的 Schrödinger 模型中，$\operatorname{Sp}(V_0)$ 的作用由线性换元给出，而 $O(U_0)$ 的作用则不再如此。两种情形都归结为描述 $\omega_\chi$ 限制到 $\operatorname{Sp}(V_0)\cdot O(U_0)$ 后的分解。

在 $\R$ 上，第1节的结果可作如下推广，参见 [Ho 2]。令 $\overline{\operatorname{Sp}}_0$ 和 $\overline O_0$ 表示 $\operatorname{Sp}(V_0)$ 与 $O(U_0)$ 在 $\overline{\operatorname{Sp}}(V)$ 中的逆像。那么，在 $\omega_\chi$ 的空间上的有界算子代数中，$\omega_\chi(\overline{\operatorname{Sp}}_0)$ 与 $\omega_\chi(\overline O_0)$ 生成彼此的交换子代数（commutant），这里采用 von Neumann 代数的意义。此外，有离散直和分解\footnote{\textbf{译者注（本次校订）：}以下离散直和分解以及 $\sigma_i$ 有限维、$\tau_i$ 为全纯型的陈述，需要 $O(U_0)$ 紧，即这里的实二次型为定号型。对于一般不定号型，应考虑直积分分解；第5.B节即讨论这一情形。}
\[
 \omega|_{\overline{\operatorname{Sp}}_0\cdot\overline O_0}
 \simeq\sum_i\tau_i\otimes\sigma_i,
\]
其中 $\tau_i$ 是 $\overline{\operatorname{Sp}}_0$ 的全纯型不可约表示，$\sigma_i$ 是 $\overline O_0$ 的不可约有限维表示；$\sigma_i$ 与 $\tau_i$ 互相决定，所得双射
\begin{equation}
 \sigma_i\longleftrightarrow\tau_i\tag{3.1}
\end{equation}
可用经典不变量理论（classical invariant theory）计算。

这些结果适用于下列对偶对：
\begin{enumerate}
\item[(i)] $\operatorname{Sp}(2nm,\R)$ 中的 $(\operatorname{Sp}(2m,\R),O(n,\R))$，其中 $n,m>0$；
\end{enumerate}
并且还可推广到：
\begin{enumerate}
\item[(ii)] $\operatorname{Sp}(2n(p+q),\R)$ 中的 $(U_n,U_{p,q})$；
\item[(iii)] $\operatorname{Sp}(4np,\R)$ 中的 $(\operatorname{Sp}(p,\mathbb O),O^*(2n))$。\footnote{\textbf{译者注（本次校订）：}此处的紧辛群是四元数酉群，通常记作 $\operatorname{Sp}(p)$ 或 $\operatorname{Sp}(p,\mathbb H)$；所用除代数应为四元数代数 $\mathbb H$。}
\end{enumerate}
因此，除推广 [Ge 1] 外，这些结果也包含了 [GK] 的早期工作（其中处理了 $n\ge2m$ 的 (i)、$n\ge2p$ 且 $p=q$ 的 (ii)，以及 $p\ge2n$ 的 (iii)），还包含了 [KV]（其中不加限制地处理 (i) 和 (ii)）。

我们将 $(\operatorname{Sp},O)$ 的上述“对偶性”所给出的对应称为对偶对应（duality correspondence）。

现在考虑任意局部域 $F$ 和 $\overline{\operatorname{Sp}}$ 的任意子群 $H$。令 $R(H)$ 表示满足下述条件的 $H$ 的不可约光滑表示（smooth representation）的等价类集合：从 $\omega_\chi$ 的光滑向量空间到该表示，存在非零的 $H$ 交织映射（intertwining map）。那么，若 $(G,G')$ 是 $\operatorname{Sp}$ 中的约化对偶对，$R(G\cdot G')$ 应当是 $R(G)$ 与 $R(G')$ 之间的一个双射（即对偶对应）的图；参见 (3.1) 和 [Ho 1]。这一猜想的一般情形尚未得到证明，但当 $G$ 或 $G'$ 紧时已经成立。后文的例子将进一步说明这一理论的由来。

\origsection{4}{阿代尔化（Adelization）}
现在令 $F$ 表示特征不为 $2$ 的 $A$ 域，$v$ 为 $F$ 的任意素位（place），$O_v$ 为 $F_v$ 的整数环，$A$ 为 $F$ 的阿代尔环（adele ring），$\chi=\prod\chi_v$ 为 $F\backslash A$ 的非平凡特征。若 $V$ 为定义在 $F$ 上的辛空间，则对每个 $v$，群及表示
\[
 H_v=H(V/F_v),\quad\operatorname{Sp}_v=\operatorname{Sp}(V/F_v),
 \quad\overline{\operatorname{Sp}}_v,\quad\rho_{\chi_v},\quad\omega_{\chi_v}
\]
均按第2节定义。

如 [Ho 1] 所述，可以定义限制张量积 $\bigotimes\omega_{\chi_v}$，由此得到 $\operatorname{Sp}_A=\operatorname{Sp}(V_A)$ 的一个乘子表示 $\omega_\chi$。若 $\overline{\operatorname{Sp}}_A$ 为由定义各 $\overline{\operatorname{Sp}}_v$ 的二次上闭链（cocycle）之积所确定的 $\operatorname{Sp}_A$ 的二重覆盖，则 $\omega_\chi$ 也定义了 $\overline{\operatorname{Sp}}_A$ 的一个普通表示。

整体表示 $\omega_\chi$ 的一个基本性质是，它在有理点 $\operatorname{Sp}(F)=\operatorname{Sp}(V/F)$ 上分裂；参见 [We]。若 $(X,Y)$ 为 $V$ 的完全极化，$\theta$ 为 Schwartz--Bruhat 空间（Schwartz--Bruhat space）$\mathcal S(X(A))$ 上由
\[
 \theta(\Phi)=\sum_{\xi\in X(F)}\Phi(\xi)
\]
定义的分布（distribution），则 $\theta$ 在 $\operatorname{Sp}(F)$ 下不变。特别地，函数
\begin{equation}
\begin{split}
 \theta_\Phi(g)&=\theta(\omega_\chi(g)\Phi)\\
 &=\sum_{\xi\in X(F)}(\omega_\chi(g)\Phi)(\xi)
\end{split}\tag{4.1}
\end{equation}
是 $\overline{\operatorname{Sp}}_A$ 上的自守函数，即增长缓慢、连续且左 $\operatorname{Sp}(F)$ 不变的函数。

振子表示整体理论的基本目标，是描述由 $\theta$ 分布通过式 (4.1) 构造出的自守形式与表示。与局部理论一样，$\operatorname{Sp}_A$ 中的约化对偶对为这一研究提供了丰富的结构。若 $(G,G')$ 是这样一对，局部理论提示，$G$ 与 $G'$ 中同 $\omega_\chi$ 交织的表示之间存在一个对偶对应。此外，这个双射应当把 $G$ 的自守表示与 $G'$ 的自守表示配对。一般情形的准确猜想见 [Ho 1]；第6节将给出具体例子。

\origsection{5}{局部例子（Local examples）}
本节始终假设 $F$ 是特征不为 $2$ 的局部域，$\chi$ 是 $F$ 的固定非平凡特征，$V_0$ 为带有反对称型 $x_1y_2-x_2y_1$ 的空间 $F^2$，$U_0$ 为带有二次型 $q$ 的正交空间 $F^n$，并令 $V=V_0\otimes U_0$。那么
\[
 \operatorname{Sp}(V_0)=\operatorname{SL}(2,F),\qquad O(U_0)=O(q),
\]
问题就是描述 $\omega_\chi$ 限制到 $\operatorname{SL}(2,F)\cdot O(U_0)$ 的情形。回忆，适当选择 $V$ 的极化后，$\omega_\chi$ 按照公式
\begin{equation}
 \omega_\chi\!\begin{pmatrix}1&b\\0&1\end{pmatrix}f(X)
 =\chi(bq(X))f(X)\tag{5.1}
\end{equation}
和
\begin{equation}
 \omega_\chi\!\begin{pmatrix}0&1\\-1&0\end{pmatrix}f(X)
 =\gamma(q,\chi)\widehat f(X)\tag{5.2}
\end{equation}
作用于 $L^2(F^n)$。正交群通过它在 $F^n$ 上的自然左作用，线性地作用于 $L^2(F^n)$。

\origsubsection{A}{各向异性情形（The anisotropic case）}
假设 $F$ 为非阿基米德局部域，$q$ 为各向异性二次型。那么 $n\le4$，$O(q)$ 紧，并且
\begin{equation}
 \omega_\chi|_{\operatorname{SL}(2,F)\cdot O(q)}
 =\sum\pi(\sigma)\otimes\sigma.\tag{5.3}
\end{equation}
这里对 $R(O(q))$ 中的 $\sigma$ 求和，每个 $\pi(\sigma)\otimes\sigma$ 的重数均为 $1$（参见 [RS 1]）。接下来需要确定 $R(O(q))$、$R(\operatorname{SL}(2,F))$ 和所得的对偶对应。

\paragraph*{情形 (i)：$n=1$.}
此时 $q(x)=ax^2$，其中 $a\in F^\times$，$O(q)=\{\pm1\}$，而相应的 $\overline{\operatorname{SL}}(2,F)$ 的表示作用于 $L^2(F)$ 的偶函数空间或奇函数空间。具体来说，若 $\sigma$ 是 $O(q)$ 的平凡表示，则 $\pi(\sigma)$（由 (5.1)、(5.2) 限制到偶函数空间定义）是 $\overline{\operatorname{SL}}(2,F)$ 的一个适当的非酉主系列表示（nonunitary principal series representation）在 $s=-1/2$ 处的唯一子表示；若 $\sigma$ 是 $O(q)$ 的非平凡表示，则 $\pi(\sigma)$（由 (5.1)、(5.2) 在奇函数空间上定义）是 $\overline{\operatorname{SL}}(2,F)$ 的超尖点表示（supercuspidal representation），并且在第6节稍后解释的意义下是“例外的”（exceptional）。详见 [Ge 2]、[G-PS] 和 [Ho 3]。注意，$ax^2$ 与 $bx^2$ 导致等价的振子表示，当且仅当 $ab^{-1}\in(F^\times)^2$。当 $F=\R$ 时，$\omega_\chi$ 的两个不可约分量平方可积（square-integrable），其极端权分别为 $1/2$ 和 $3/2$。\footnote{\textbf{译者注（本次校订）：}这里关于平方可积性的陈述须加修正：最低权为 $1/2$ 的偶分量不平方可积，最低权为 $3/2$ 的奇分量才属于离散系列。实 metaplectic 群的最低权表示平方可积所需的条件为最低权大于 $1$；取相反的加性特征时，应相应改说最高权及其绝对值。}

\paragraph*{情形 (ii)：$n=2$.}
令 $U_0$ 为 $F$ 的二次扩张 $K$，其内积由范数型（norm form）$q$ 导出。那么 $O(q)$ 是 $K$ 中的范数为 $1$ 的群 $K^1$ 与 $K/F$ 的 Galois 群的半直积。此时，$R(O(q))$ 中的每个 $\sigma$ 都可通过 $K^1$ 的特征来描述，而 $R(\operatorname{SL}(2,F))$ 的大部分表示都是超尖点的。详见 [ST]、[Cas]、[G] 或 [RS 1]；在 [ST] 中，由于用特殊正交群（special orthogonal group）$K^1$ 代替 $O(U_0)$，对应是二对一的。这个例子具有重要的历史意义：Shalika 和 Tanaka 正是借助振子表示，首次构造出了这样一类不可约表示。

\paragraph*{情形 (iii)：$n=3$ [RS 1].}
令 $H$ 为 $F$ 上唯一的除四元数代数（division quaternion algebra），$H^0$ 为纯四元数的子空间，$q_3$ 为 $H$ 上约化范数型（reduced norm form）在 $H^0$ 上的限制。那么，$F$ 上任意各向异性三元型 $q$ 都等价于某个 $aq_3$，其中 $a\in F^\times/(F^\times)^2$；$O(q)$ 的每个不可约表示 $\sigma$ 都属于 $R(O(q))$（$\operatorname{SO}(q)$ 同构于 $(H^0)^\times$）。\footnote{\textbf{译者注（本次校订）：}此处有两点需修正。首先，纯四元数空间 $H^0$ 对乘法不封闭，正确的群同构为 $\operatorname{SO}(q)\simeq H^\times/F^\times$，由共轭作用给出。其次，并非 $O(q)$ 的每个不可约表示都出现在这里：行列式特征便不出现。事实上，每个向量的稳定子都含一个行列式为 $-1$ 的反射，因此按行列式特征变换的函数只能恒为零。}[RS 1] 还证明，相应的 $\overline{\operatorname{SL}}(2,F)$ 的表示 $\pi(\sigma)$ 平方可积，并且在 $F$ 的剩余特征为奇数时，所得表示穷尽全部平方可积的“真正的”表示。

\paragraph*{情形 (iv)：$n=4$.}
令 $U_0$ 表示定义在 $F$ 上的除四元数代数 $H$，其内积由约化范数型 $q$ 导出。那么，虽然 $O(q)$ 与 $H^\times$ 不同，但 $R(O(q))$ 中的每个 $\sigma$ 自然地对应于 $H^\times$ 的一个不可约表示，$H^\times$ 的所有此类表示都由此产生，并且 $R(\operatorname{SL}(2,F))$ 中的每个 $\pi(\sigma_v)$ 都平方可积。事实上，$\pi(\sigma_v)$ 唯一地扩张为 $\operatorname{GL}(2,F_v)$ 的一个表示 $\pi'(\sigma_v)$\footnote{\textbf{译者注（本次校订）：}仅凭对 $\operatorname{SL}(2,F_v)$ 的限制和这里规定的中心特征，不能保证扩张的唯一性：用 $\eta\circ\det$ 扭转，其中 $\eta$ 为二次特征，既不改变限制，也不改变中心特征。例如 Steinberg 表示与其非平凡二次扭转具有这两项相同的数据。此处的特定扩张须由所述对应本身确定。}，使
\[
 \pi'(\sigma_v)\!\begin{pmatrix}a&0\\0&a\end{pmatrix}=\sigma_v(a)I,
\]
所得对偶对应 $\sigma_v\leftrightarrow\pi^{\prime}(\sigma_v)$ 满射到 $\operatorname{GL}(2,F_v)$ 的平方可积不可约表示的等价类集合；详见 [JL]。

注意，情形 (iv) 与情形 (iii) 相似：局部对偶对应中的 $\pi(\sigma)$ 由其平方可积性刻画。只有情形 (ii) 尚缺少对 $R(O(q))$ 的简洁刻画。

\origsubsection{B}{非紧（实）情形（The noncompact real case）}
假设 $F=\R$，$\chi(u)=e^{2\pi iu}$。令 $O(q)=O(k,l)$ 表示 $\R^n$ 上符号（signature）为 $(k,l)$ 的标准二次型的等距群，其中 $n=k+l$。Strichartz 在 [St] 中已基本给出 $O(q)$ 在 $L^2(\R^n)$ 中的自然作用的分解，也就是 $\omega_\chi$ 在 $O(q)$ 上的限制的分解。利用 [Re] 关于 $\operatorname{SL}(2,\R)$ 表示的张量积的结果，Howe 证明了他的局部对偶猜想的如下版本：
\[
 \left.\omega_\chi\right|_{\overline{\operatorname{SL}}(2,\R)\cdot O(q)}
 \simeq\int\sigma_s'\otimes\sigma_s\,ds,
\]
其中 $ds$ 是 $\operatorname{SL}(2,\R)$ 的酉对偶（unitary dual）上的一个 Borel 测度，$\sigma_s$ 和 $\sigma_s'$ 分别为 $O(k,l)$ 和 $\overline{\operatorname{SL}}(2,\R)$ 的不可约酉表示，并且对几乎所有 $s$，二者互相决定。所得对应 $\sigma_s\mapsto\pi(\sigma_s)=\sigma_s'$，对于 $R(\overline{\operatorname{SL}}(2,\R))$ 和 $R(O(q))$ 中的离散系列表示尤其有趣。完整细节见 [Ho 4] 和 [RS 2]；情形 $(k,l)=(2,1)$ 在 [Ge 2] 中有所描述。

\origsubsection{C}{非紧 p 进情形（Noncompact p-adic case）}
把例 5.A(ii) 的二次扩张替换为带内积
\[
 ((x,y),(x',y'))=xx'-yy'
\]
的双曲平面（hyperbolic plane）$F^2$，就得到一个有趣的例子。此时 $O(U_0)$ 分裂。因此，可以找到 $V$ 的一个被 $\operatorname{Sp}(V_0)=\operatorname{SL}(2,F)$ 保持的极化，取
\[
 X_0=\{(x,x):x\in F\},\qquad Y_0=\{(x,-x):x\in F\}.
\]
在相应的 Schrödinger 模型中，$\operatorname{SL}(2,F)$ 通过线性换元作用于 $L^2(F^2)$。具体地，
\[
 \omega_\chi(g)f\!\begin{pmatrix}x\\y\end{pmatrix}
 =f\!\left(g^{-1}\!\begin{pmatrix}x\\y\end{pmatrix}\right).
\]
由于 $F^2-\{0\}$ 上的函数可与 $G$ 上在
\[
 N=\left\{\begin{pmatrix}1&b\\0&1\end{pmatrix}\right\}
\]
下不变的函数等同，振子表示在此情形下给出 $\operatorname{SL}(2,F)$ 主系列表示的直积分（direct integral）。此时，正交群的作用不再仅由线性换元给出。事实上，对应于反射（把 $(x,y)$ 送到 $(y,x)$ 的元素）的算子，是 $L^2(F^2)$ 中相对于 $F^2$ 上交替型的 Fourier 变换。\footnote{\textbf{译者注（本次校订）：}交换 $(x,y)$ 与 $(y,x)$ 只有在双曲坐标中才表示这里的反射。在本节起初采用的坐标中，二次型为 $x^2-y^2$，交换坐标会改变其符号；交换两条各向同性直线 $X_0,Y_0$ 的等距反射应写作 $(x,y)\mapsto(x,-y)$。}

Cartier 的讲座已讨论这些事实在 $\operatorname{SL}(2,F)$ 表示论中的意义（亦见 [G]）。

\origsection{6}{整体例子（Global examples）}
\origsubsection{A}{各向异性情形（The anisotropic case）}
\paragraph*{情形 (i)：$O(1)$.}
将 $L^2(F_v)$ 中振子表示的各个分量作限制张量积，可构造出一批 $\overline{\operatorname{SL}}(2)$ 的自守表示，并通过诱导得到 $\overline{\operatorname{GL}}(2)$ 的自守表示。Howe 最初向我指出，振子表示的非平凡分量可以用来构造尖点形式（参见 [Ho 3]）。下面简述这一思路的后续发展，细节见 [G-PS]。

若 $\sigma=\bigotimes\sigma_v$ 为 $O(1)$ 的自守表示，则除某个基数为偶数的有限集合 $S$ 中的素位外，$\sigma_v$ 对所有 $v$ 都是平凡表示。若 $\chi=\prod\chi_v$ 为 $A/F$ 的特征，令 $\omega_{\chi_v}$ 表示 $\overline{\operatorname{SL}}(2,F_v)$ 在 $L^2(F_v)$ 中与 $\chi_v$ 对应的振子表示。固定 $A^\times/F^\times$ 的一个特征 $\lambda=\prod\lambda_v$，使 $v\notin S$ 时 $\lambda_v(-1)=1$，$v\in S$ 时 $\lambda_v(-1)=-1$；相应地，以 $\omega_{\chi_v}'$ 表示 $\omega_{\chi_v}$ 的偶部分或奇部分。再将 $\omega_{\chi_v}'$ 扩张到
\[
 G_v^*=\{g\in\operatorname{GL}(2,F_v):\det(g)\in(F_v^\times)^2\},
\]
方法是把 $\omega_{\chi_v}'\!\left(\begin{smallmatrix}1&0\\0&a^2\end{smallmatrix}\right)$ 定义为算子
\[
 \Phi(x)\longmapsto\lambda(a)|a|^{-1/2}\Phi(a^{-1}x),
\]
再把 $\omega_{\chi_v}'$ 诱导到 $\operatorname{GL}(2,F_v)$，得到 $\operatorname{GL}(2,F_v)$ 的二重覆盖 $\overline G_v$ 的表示 $\omega_v'$。这个 $\omega_v'$ 是 $\overline G_v$ 的不可约、酉的“真正的”表示；对所有 $v$，它都与 $\chi_v$ 的选择无关，而且对几乎所有 $v$，它都属于第一类（class 1）。

表示 $\pi(\lambda)=\bigotimes\omega_v'$ 是 $\overline{\operatorname{GL}}(2,A)$ 的一个真正的自守表示，而且当且仅当 $S\ne\varnothing$ 时为尖点表示。一般地，当 $S$ 为空时，这些 $\pi(\lambda)$ 生成 $\overline{\operatorname{GL}}(2,A)$ 的离散非尖点谱（discrete noncuspidal spectrum）（参见 [G-S]）；特别地，当 $F=\Q$ 时，$\pi(\lambda)$ 推广了经典 theta 函数
\[
 \theta(z)=\sum_{n=-\infty}^{\infty}\exp(2\pi i n^2z).
\]
另一方面，当 $S\ne\varnothing$ 时，$\pi(\lambda)$ 推广了 Dedekind 的 $\eta$ 函数等经典尖点形式。此时，$\pi(\lambda)$ 是例外的，意思是其表示空间中任意函数的 Fourier 展开“只包含一个特征轨道”；这推广了 $\eta$ 函数展开“只有平方项”这一事实：\footnote{\textbf{译者注（本次校订）：}若 $(3/n)$ 按通常的 Kronecker 符号理解，下式应改用 $(12/n)$，或者将求和限制在奇数 $n$。例如 $\eta(24z)=q-q^{25}-q^{49}+\cdots$（$q=e^{2\pi iz}$）没有 $q^4$ 项，而 $(3/2)=-1$。}
\[
 \eta(24z)=\sum_{n=1}^\infty(3/n)\exp(2\pi i n^2z).
\]
由于 $\overline{\operatorname{GL}}(2,A)$ 上尖点形式 $\bar\pi$ 的 Fourier 展开，是相对于同构于 $F$ 与 $F^\times$ 的二重覆盖 $\overline F^\times$ 的乘积的一个群进行的，我们可以这样说：若给定 $F$ 上的 $\chi$ 后，$\bar\pi$ 相对于 $(\chi,\bar\mu)$ 的 Whittaker 模型（Whittaker model）只对 $\overline F^\times$ 的一个特征 $\bar\mu$ 存在，则称 $\bar\pi$ 为例外的。因此，例外性的概念在局部也有意义，详见 [G-PS]。我们的猜想是，振子表示在局部和整体上都产生所有可能的例外表示。换言之，对于 $O(1)$，对偶对应 $\sigma\mapsto\pi(\sigma)$ 的像应是 $\overline G$ 的所有例外真正表示。

刚才描述的猜想的整体部分，用经典语言表述为：任何权为 $k/2$ 且 Fourier 展开只含平方项的尖点形式，其权必为 $1/2$ 或 $3/2$，并且是如下形式的函数的平移的线性组合：
\[
 \sum_{n=1}^{\infty}\psi(n)n^\nu\exp(2\pi i n^2z),\qquad\nu=0\text{ 或 }1.
\]
我们在本次会议上得知 M.-F. Vignéras 已在 [V] 中证明这一断言。

\paragraph*{情形 (ii)：$O(2)$.}
令 $K$ 为域 $F$ 的二次扩张，$\chi$ 为 $F\backslash A$ 的一个特征，$\sigma=\bigotimes\sigma_v$ 为 $K^1(F)\backslash K^1(A)$ 的一个特征。将第5.A(ii)节构造的局部表示 $\pi(\sigma_v)$（$\operatorname{SL}(2,F_v)$ 的表示）作限制张量积，便得到 $\operatorname{SL}(2,A)$ 的自守尖点表示，只要 $\sigma$ 非平凡即可。这些自守形式与 E. Hecke [H] 和 H. Maass [Ma] 较早构造的带广义 Hecke 特征的 theta 级数有关；详见 [ST] 和 [Ge 3，第7章]。

\paragraph*{情形 (iii)：$O(4)$.}
令 $D$ 为 $F$ 上的除四元数代数，$S_D$ 为 $D$ 分歧（ramify）的素位集合。设 $\sigma=\bigotimes\sigma_v$ 属于 $R(D^\times(A))$。若 $v\in S_D$，则 $D_v$ 是 $F_v$ 上唯一的除四元数代数，并如第5节那样定义 $\pi'(\sigma_v)$。另一方面，当 $v\notin S_D$ 时，$D_v^\times$ 同构于 $\operatorname{GL}(2,F_v)$，自然对偶对应就是恒等映射。

现在进一步假设 $\sigma$ 是自守表示，即 $\sigma$ 可以嵌入 $G'(A)=D^\times(A)$ 上的自守形式空间。那么，按 [JL]，$\operatorname{GL}(2,A)$ 的表示 $\pi'(\sigma)=\bigotimes_v\pi'(\sigma_v)$ 也是自守的。反过来，若 $\pi=\bigotimes\pi_v$ 属于 $R(\operatorname{GL}(2,A))$，也就是说，对某些 $\sigma_v\in R(D_v^\times)$，有 $\pi_v=\pi'(\sigma_v)$（等价地，对每个 $v\in S_D$，$\pi_v$ 都平方可积），那么只有当 $\sigma=\bigotimes\sigma_v$ 自守时，$\pi$ 才自守；参见 [JL] 的定理14.4和16.1。这些结果在除四元数代数这一特例中证实了 Howe 理论的主要整体猜想。[JL] 的证明使用迹公式（trace formula），Howe 也已得到一个独立证明，所用方法不同，而且更为初等。

\origsubsection{B}{$O(2,1)$}
令 $U_0$ 为 $F$ 上迹为零的 $2\times2$ 矩阵构成的三维空间，内积由行列式函数导出。置 $V_0=(F^2,\langle\ ,\ \rangle)$，$V=V_0\otimes U_0$。那么对 $F$ 的每个素位 $v$，$(\operatorname{SL}(2,F_v),O(U_0))$ 是 $\operatorname{Sp}(V)$ 中的约化对偶对。

对 $F_v=\R$，对偶对应 $\sigma_v\leftrightarrow\pi(\sigma_v)$ 把 $\operatorname{PGL}(2,\R)=\operatorname{SO}(2,\R)$ 的权为 $k-1$ 的离散系列表示，与 $\overline{\operatorname{SL}}(2,\R)$ 的权为 $k/2$ 的离散系列表示配对。\footnote{译注：原书此处确实印作 $\operatorname{SO}(2,\R)$；结合本段 $O(2,1)$ 的主题，该记号疑有原印刷遗漏。这里保留原记号。} 当 $F=\Q$ 时，[Ge 2] 猜想所得整体对应应当推广 Shimura [Shm] 在权为 $k/2$ 和 $k-1$ 的经典形式之间的对应。Howe 最近对第一类表示的局部对偶对应的描述（[Ho 1]），以及 Shintani 和 Niwa 利用 theta 函数积分研究整体对应的工作（[Sht] 和 [N]），为这一点提供了证据。关于此问题的进一步工作，亦见 [G-PS]。

若 $F_v$ 为非阿基米德局部域，$\sigma_v$ 是 $\operatorname{PGL}(2,F_v)$ 的一个形如 $\pi'(\sigma_v')$ 的离散系列表示，其中 $\sigma_v'$ 是一个各向异性三元型的正交群的适当有限维表示，那么 $\pi(\sigma_v)$ 应当与第5.A节情形 (iii) 中描述的 $\overline{\operatorname{SL}}(2,F_v)$ 的表示 $\pi(\sigma_v')$ 相同。

\origsubsection{C}{$O(2,2)$（$F=\Q$）}
令 $K=\Q(\sqrt\Delta)$，其中 $\Delta>0$ 为判别式（discriminant），以 $\tau$ 表示 $K/\Q$ 的 Galois 自同构。令
\[
 U_0=\left\{X=\begin{pmatrix}x_1&x_4\\x_3&-x_1^\tau\end{pmatrix}
 :x_1\in K,\ x_3,x_4\in\Q\right\}\simeq\Q^4,
\]
并由 $q(X)=-2\det(X)$ 定义 $q:U_0\to\Q$。那么 $q$ 的符号为 $(2,2)$，相应的整体对偶对应应当把整权模形式（modular form of integral weight）与 Hilbert 模形式（Hilbert modular form）配对，因为存在从 $\operatorname{SL}(2,\R)\times\operatorname{SL}(2,\R)$ 到 $\operatorname{SO}(2,2)$ 的自然同态。已有下列作者用上半平面上形式的经典语言研究了这一对应：Oda [O]、Asai [A]、Kudla [K]、Rallis 和 Schiffmann [RS 3]、Zagier [Z]；Doi--Naganuma [DN] 则更早从稍有不同的角度开展了相关研究。

\paragraph*{校样补记（Note added in proof）.}
我们现已证明原第293页下部描述的整体“例外性猜想”。局部猜想已经由 James Meister 证明，将刊载于他即将完成的 Cornell 大学博士论文中。

\begin{thebibliography}{G-PS}
\bibitem[A]{c18-A} T. Asai, \emph{On the Doi-Naganuma lifting associated to imaginary quadratic fields}, 即将发表。
\bibitem[Car]{c18-Car} P. Cartier, \emph{Representations of $p$-adic groups: A survey}, 本论文集，第1册，111--155页。
\bibitem[Cas]{c18-Cas} W. Casselman, \emph{On the representations of $\operatorname{SL}_2(k)$ related to binary quadratic forms}, Amer. J. Math. \textbf{94} (1972), 810--834.
\bibitem[DN]{c18-DN} K. Doi and H. Naganuma, \emph{On the functional equation of certain Dirichlet series}, Invent. Math. \textbf{9} (1969), 1--14；另见 H. Naganuma, \emph{On the coincidence of two Dirichlet series associated with cusp forms of Hecke's ``Neben''-type and Hilbert modular forms over a real quadratic field}, J. Math. Soc. Japan \textbf{25} (1973), 547--554.
\bibitem[Ge 1]{c18-Ge1} S. Gelbart, \emph{Holomorphic discrete series for the real symplectic group}, Invent. Math. \textbf{19} (1973), 49--58.
\bibitem[Ge 2]{c18-Ge2} S. Gelbart, \emph{Weil's representation and the spectrum of the metaplectic group}, Lecture Notes in Math., vol. 530, Springer-Verlag, Berlin, 1976.
\bibitem[Ge 3]{c18-Ge3} S. Gelbart, \emph{Automorphic forms on adele groups}, Ann. of Math. Studies, no. 83, Princeton Univ. Press, Princeton, N.J., 1975.
\bibitem[G-PS]{c18-GPS} S. Gelbart and I. I. Piatetski-Shapiro, \emph{Automorphic $L$-functions of half-integral weight}, Proc. Nat. Acad. Sci. U.S.A. \textbf{75} (1978), 1620--1623.
\bibitem[G-S]{c18-GS} S. Gelbart and P. J. Sally, \emph{Intertwining operators and automorphic forms for the metaplectic group}, Proc. Nat. Acad. Sci. U.S.A. \textbf{72} (1975), 1406--1410.
\bibitem[G]{c18-G} P. Gérardin, \emph{Représentations des groupe $\operatorname{SL}_2$ d'un corps local (d'après Gelfand, Graev, et Tanaka)}, Séminaire Bourbaki, $20^e$ année, 1967/1968, no. 322, Benjamin, New York and Amsterdam, 1969.
\bibitem[GK]{c18-GK} K. Gross and R. Kunze, \emph{Bessel functions and representation theory. II}, J. Functional Analysis \textbf{25} (1977), 1--49.
\bibitem[H]{c18-H} E. Hecke, \emph{Mathematische Werke}, Vandenhoeck, Göttingen, 1959.
\bibitem[Ho 1]{c18-Ho1} R. Howe, \emph{$\theta$-series and invariant theory}, 本论文集，第1册，275--285页。
\bibitem[Ho 2]{c18-Ho2} R. Howe, \emph{Appendix to remarks on classical invariant theory}, 预印本，1976。
\bibitem[Ho 3]{c18-Ho3} R. Howe, \emph{Review of [GS]}, Zentralblatt für Math., Band 303, Review 22010, January 1976, p. 98.
\bibitem[Ho 4]{c18-Ho4} R. Howe, \emph{On some results of Strichartz and of Rallis and Schiffmann}, 预印本，1976。
\bibitem[Ho 5]{c18-Ho5} R. Howe, \emph{Invariant theory and duality for classical groups over finite fields, with applications to their singular representation theory}, 预印本。
\bibitem[JL]{c18-JL} Hervé Jacquet and Robert Langlands, \emph{Automorphic forms on GL(2)}, Lecture Notes in Math., vol. 114, Springer-Verlag, Berlin, 1970.
\bibitem[KV]{c18-KV} M. Kashiwara and M. Vergne, \emph{On the Segal-Shale-Weil representation and harmonic polynomials}, 预印本，1976。
\bibitem[K]{c18-K} S. Kudla, \emph{Relations between automorphic forms produced by theta-functions}, Modular Functions of One Variable. VI, Lecture Notes in Math., vol. 627, Springer-Verlag, Berlin, 1977.
\bibitem[Kz]{c18-Kz} D. Kazhdan, \emph{Applications of Weil's representation}, 预印本，1976。
\bibitem[Ma]{c18-Ma} H. Maass, \emph{Über eine neue Art von nicht analytischen automorphen Funktionen und die Bestimmung Dirichletscher Reihen durch ihre Funktionalgleichung}, Math. Ann. \textbf{121} (1949), 141--183.
\bibitem[Mo]{c18-Mo} Calvin Moore, \emph{Group extensions of $p$-adic linear groups}, Inst. Hautes Études Sci. Publ. Math. \textbf{35} (1968), 157--222.
\bibitem[N]{c18-N} S. Niwa, \emph{Modular forms of half-integral weight and the integral of certain theta-functions}, Nagoya Math. J. \textbf{56} (1975), 147--161.
\bibitem[O]{c18-O} T. Oda, \emph{On modular forms associated with indefinite quadratic forms of signature $(2,n-2)$}, 预印本。
\bibitem[RS 1]{c18-RS1} S. Rallis and G. Schiffmann, \emph{Représentations supercuspidales du groupe metaplectique}, 预印本，1977。
\bibitem[RS 2]{c18-RS2} S. Rallis and G. Schiffmann, \emph{Discrete spectrum of the Weil representation}, 预印本，1975。
\bibitem[RS 3]{c18-RS3} S. Rallis and G. Schiffmann, \emph{Automorphic forms constructed from the Weil representation, holomorphic case}, 预印本，1976。
\bibitem[Re]{c18-Re} J. Repka, \emph{Tensor products of unitary representations of $\operatorname{SL}_2(\R)$}, 预印本。
\bibitem[Sa]{c18-Sa} M. Saito, \emph{Représentations unitaires des groups symplectiques}, J. Math. Soc. Japan \textbf{24} (1972), 232--251.
\bibitem[ST]{c18-ST} Joseph Shalika and S. Tanaka, \emph{On an explicit construction of a certain class of automorphic forms}, Amer. J. Math. \textbf{91} (1969), 1049--1076.
\bibitem[Shz]{c18-Shz} H. Shimizu, \emph{Theta-series and automorphic forms on GL(2)}, J. Math. Soc. Japan \textbf{24} (1972), 638--683.
\bibitem[Shm]{c18-Shm} G. Shimura, \emph{On modular forms of half-integral weight}, Ann. of Math. (2) \textbf{97} (1973), 440--481.
\bibitem[Sht]{c18-Sht} T. Shintani, \emph{On construction of holomorphic cusp forms of half-integral weight}, Nagoya Math. J. \textbf{58} (1975), 83--126.
\bibitem[Sie]{c18-Sie} C. L. Siegel, \emph{Indefinite quadratische Formen und Funktionen Theorie. I, II}, Math. Ann. \textbf{124} (1951/52), 17--54, 364--387.
\bibitem[St]{c18-St} R. Strichartz, \emph{Harmonic analysis on hyperboloids}, J. Functional Analysis \textbf{12} (1973), 341--383.
\bibitem[V]{c18-V} M.-F. Vignéras, \emph{Facteurs gamma et équations fonctionnelles}, Modular Functions of One Variable. VI, Lecture Notes in Math., vol. 627, Springer-Verlag, Berlin, 1977.
\bibitem[We]{c18-We} André Weil, \emph{Sur certaines groupes d'opérateurs unitaires}, Acta Math. \textbf{111} (1964), 143--211.
\bibitem[Z]{c18-Z} D. Zagier, \emph{Modular forms associated to real quadratic fields}, Invent. Math. \textbf{30} (1975), 1--46.
\end{thebibliography}
\noindent Cornell 大学（Cornell University）
